% Generated draft: requires mathematical and source review.
\documentclass[11pt]{article}
\usepackage[margin=25mm]{geometry}
\usepackage{fontspec}
\setmainfont{Latin Modern Roman}
\usepackage{amsmath,amssymb,mathtools,unicode-math}
\setmathfont{Latin Modern Math}
\usepackage{enumitem,needspace,xurl,hyperref}
\setlength{\parindent}{0pt}
\setlength{\parskip}{5pt}
\setlength{\emergencystretch}{3em}
\newcommand{\N}{\mathbb{N}}
\newcommand{\Z}{\mathbb{Z}}
\newcommand{\Q}{\mathbb{Q}}
\newcommand{\R}{\mathbb{R}}
\newcommand{\C}{\mathbb{C}}
\newcommand{\F}{\mathbb{F}}
\newcommand{\A}{\mathbb{A}}
\newcommand{\PP}{\mathbb{P}}
\newcommand{\E}{\mathbb{E}}
\newcommand{\eps}{\varepsilon}
\newcommand{\dd}{\,\mathrm{d}}
\begin{document}
\section*{UnsolvedMath ID: 30005159}
\textbf{Source problem code:} catalogue entry
\subsection{Definitions and mathematical statement}
The supplied registry asks whether, for every integer \(d\geq4\), the function \(f_\star(\tau,\xi)=e^{-|\tau|}\) is a global maximizer of the \(p=2\) Fourier adjoint-restriction inequality on \(K_d^+=\{(\tau,\xi):\tau=|\xi|\}\). However, the excerpt omits the referenced inequality (2), the ambient dimension convention, the cone measure, and the target Lebesgue exponent. Thus the following precise formulation is conditional on those missing data.

Let \(n_d\) denote the intended spatial dimension, so that
\[
K_d^+=\{(\tau,\xi)\in\R\times\R^{n_d}:\tau=|\xi|\}.
\]
Let \(\sigma_d\) be the intended measure on this cone, and let \(q_d\) be the target Lebesgue exponent in inequality (2) when \(p=2\). With one Fourier normalization, the associated extension operator is initially
\[
(E_d f)(t,x)=\int_{K_d^+} e^{i(t\tau+x\cdot\xi)}f(\tau,\xi)\,\mathrm d\sigma_d(\tau,\xi),
\qquad (t,x)\in\R\times\R^{n_d},
\]
for functions for which this integral is defined. Assume that the intended inequality supplies a bounded extension
\[
E_d:L^2(K_d^+,\sigma_d)\longrightarrow L^{q_d}(\R^{1+n_d}).
\]
Define its sharp constant by
\[
C_d=\sup_{0\ne f\in L^2(K_d^+,\sigma_d)}
\frac{\|E_d f\|_{L^{q_d}(\R^{1+n_d})}}
{\|f\|_{L^2(K_d^+,\sigma_d)}}.
\]
The question is whether, for every integer \(d\geq4\), the function
\[
f_\star(\tau,\xi)=e^{-|\tau|}=e^{-|\xi|}\quad\text{on }K_d^+
\]
is a nonzero member of the intended \(L^2\) space and satisfies
\[
\|E_d f_\star\|_{L^{q_d}(\R^{1+n_d})}
=C_d\|f_\star\|_{L^2(K_d^+,\sigma_d)}.
\]
Here \(n_d\), \(\sigma_d\), and \(q_d\) remain unspecified source data; this template does not identify them or certify a complete reconstruction of (2).

The original excerpt also reports that \(f_\star\) is a global maximizer on \(K_d^+\) for \(d=2,3\), and a local maximizer there for \(d\geq3\). Separately, it reports global maximality on a cone denoted \(K_d\) for \(d=3\), local maximality for all odd \(d\), and failure even to be a critical point for even \(d\). The excerpt does not define \(K_d\), the topology for local maximality, or the variational meaning of criticality. These are supplied background claims, not assertions independently verified here, and the claims about \(K_d\) must not be transferred to \(K_d^+\).
\subsection{Short English statement}
For each integer \(d\geq4\), determine whether the exponential profile attains the sharp \(L^2\) cone extension bound on the positive cone. The supplied excerpt omits the defining inequality and its measure, exponent, and dimension convention, so the displayed formulation remains conditional. The supplied literature assessment also mixes qualifications concerning two differently denoted cones.
\subsection{Sources}
\begin{itemize}
\item Supplied local UnsolvedMath registry record, permanent ID \(30005159\), source problem code \texttt{catalogue entry}, title \emph{Exponential Maximizers for Cone Adjoint Restriction}, in the Oberwolfach Reports Open Problems collection. No individual proposer is supplied; the proposed year is recorded as 2022.
\item Supplied source citation: \emph{Real Analysis, Harmonic Analysis and Applications} (2023), Oberwolfach Reports, pp.\ 1743--1803, DOI \texttt{10.4171/owr/2022/31}. \url{https://doi.org/10.4171/owr/2022/31}. The supplied excerpt refers to inequality (2), Theorem 2, references [2, 3, 4], and Table 1, without reproducing the inequality, bibliography entries, or table.
\item Supplied literature source: Negro, cone extension maximizers. \url{https://doi.org/10.1112/jlms.70112}.
\item Supplied literature source: Goncalves--Negro, local maximizers. \url{https://arxiv.org/abs/2003.11955}.
\item Supplied literature source: OWR report. \url{https://ems.press/content/serial-article-files/46969}.
\item Literature assessment dated exactly \texttt{2026-08-22}: supplied dated curation, not independently verified here. It labels the problem partially solved and reports global maximality in dimensions two and three, local maximality in several higher dimensions, and failure of criticality in other dimensions, concluding that the dimension-free question has no single affirmative answer. The supplied background explicitly requests independent verification of these time-sensitive claims.
\item The registry labels the statement \texttt{corrected\_verified} and says that the official report was checked and its notation repaired. This describes supplied prior curation; no source visit or independent verification is claimed in this draft.
\end{itemize}
\end{document}
